% Note technique pour l'équipe GED : versement des documents de la gestion des thèses.
% Compilation : pdflatex note-ged.tex (deux fois), depuis docs/ged.
\documentclass[11pt,a4paper]{article}

\newcommand{\titreguide}{Note technique GED}
\newcommand{\versionguide}{1.0}
\newcommand{\dateguide}{28 octobre 2026}
\newcommand{\versionthese}{0.19.0}

\usepackage{../guides/guide-commun}
\usepackage{fancyvrb}
\fvset{fontsize=\footnotesize,frame=single,framesep=2mm,rulecolor=\color{cadre}}

\begin{document}

\begin{center}
  {\LARGE\bfseries\color{moodle} Versement des documents des thèses dans la GED}\\[0.6em]
  {\large Stockage des documents scellés et extraction}\\[0.8em]
  {\small ThesisManager (\texttt{mod\_these}) \versionthese{} -- note \versionguide{} du \dateguide{} -- Sébastien Lefebvre}
\end{center}

\vspace{1em}

\begin{astuce}
\textbf{En bref.} Les documents signés sont figés (\og scellés\fg{}) dans la zone de fichiers
\texttt{sealed} du module, gérée par l'API de fichiers de Moodle. Le connecteur lit la table
\texttt{mdl\_files} et le dossier \texttt{moodledata}, \textbf{en lecture seule}, vérifie les empreintes
SHA-1 et SHA-256, et identifie chaque version par \emph{dossier + document + empreinte}. Le script
\texttt{docs/ged/extraction\_ged.py} fait tout cela et produit un manifeste CSV.
\end{astuce}

\tableofcontents

% ===========================================================================
\section{Objet et périmètre}

La gestion des thèses d'exercice vétérinaire se fait dans Moodle, avec le module ThesisManager. Chaque
étudiant a un \textbf{dossier} qui réunit son sujet, sa soutenance, son jury et ses documents. Les
documents officiels sont produits et signés dans le module.

Cette note décrit, pour l'équipe chargée de la GED :
\begin{itemize}
  \item comment un document est scellé et ce que garantit le scellé (section~\ref{sec:scellement}) ;
  \item où et comment les fichiers sont stockés (section~\ref{sec:stockage}) ;
  \item quels fichiers verser et avec quelles métadonnées (sections~\ref{sec:zones} et~\ref{sec:metadonnees}) ;
  \item comment les extraire et vérifier leur intégrité (sections~\ref{sec:acces} à~\ref{sec:script}) ;
  \item les règles que le connecteur doit respecter (section~\ref{sec:regles}).
\end{itemize}

% ===========================================================================
\section{Le scellement}\label{sec:scellement}

\paragraph{Documents concernés.} Neuf documents de l'établissement sont signés dans le module :

\begin{center}
\small
\begin{tabularx}{\linewidth}{@{}l X l@{}}
\toprule
\textbf{Fichier} & \textbf{Document} & \textbf{Signataires} \\
\midrule
\texttt{subject.pdf} & Dépôt de sujet (E01) & étudiant, directeur, DEVE \\
\texttt{bordereau.pdf} & Bordereau de thèse (ABES) & étudiant \\
\texttt{report.pdf} & Rapport du directeur & directeur \\
\texttt{grid.pdf} & Grille de conformité & directeur \\
\texttt{agreement.pdf} & Fiche d'agréments (E02) & DEVE, directeur, DG, CCEM \\
\texttt{signaletique.pdf} & Fiche signalétique (E03) & président du jury, DG, CCEM \\
\texttt{contract.pdf} & Contrat de diffusion électronique & étudiant, directeur, DG, université \\
\texttt{conformity.pdf} & Attestation de conformité du manuscrit & bibliothèque \\
\texttt{conclusions.pdf} & Conclusions de la thèse & directeur, président, DG, CCEM \\
\bottomrule
\end{tabularx}
\end{center}

La direction générale (DG) signe toujours par délégation, par un agent de la DEVE. Le directeur
extérieur éventuel signe aussi le dépôt de sujet et la fiche d'agréments.

\paragraph{Quand.} Dès que la \textbf{dernière signature attendue} d'un document est apposée, le module
produit son PDF une dernière fois et le conserve tel quel : c'est le document \textbf{scellé}. Il ne
dépend plus ensuite ni des modèles de mise en forme ni du contrat type.

\paragraph{Ce qu'il contient.} Pour chaque signataire : le nom, la date, l'image de la signature,
l'empreinte du contenu signé et, le cas échéant, la mention \og par délégation\fg{}, la vérification
d'identité par la DEVE (signature sur place) ou la signature à la main (scan joint au dossier).

\paragraph{Ce que garantit le module.}
\begin{itemize}
  \item Chaque signature conserve l'empreinte SHA-256 des données signées : si le contenu du dossier
        change ensuite, la signature est invalidée et le document doit être signé de nouveau.
  \item À la création du scellé, l'empreinte \textbf{SHA-256 du PDF}, calculée sur ses octets, est
        inscrite à l'historique du dossier (table \texttt{mdl\_these\_log}, action
        \texttt{documentsealed}, remarque \og \textit{Libellé du document} — SHA-256
        \textit{empreinte}\fg{}).
\end{itemize}

\begin{attention}
\textbf{Un scellé peut être levé.} Le fichier scellé est supprimé quand une de ses signatures disparaît :
\begin{itemize}
  \item invalidation d'une signature (données du dossier modifiées après signature) ;
  \item ajournement de la thèse (toutes les signatures sauf celles du dépôt de sujet) ;
  \item changement de sujet (signatures du dépôt de sujet ; l'ancien dépôt signé est d'abord archivé
        dans la zone \texttt{subjectarchive}) ;
  \item effacement des données d'un utilisateur au titre du RGPD.
\end{itemize}
Une fois toutes les signatures de nouveau réunies, le document est scellé à nouveau, \textbf{sous le
même nom de fichier}, avec un autre contenu, une autre empreinte et une nouvelle ligne
\texttt{documentsealed} à l'historique. La GED doit conserver les versions déjà versées.
\end{attention}

% ===========================================================================
\section{Stockage dans Moodle}\label{sec:stockage}

Les fichiers ne sont pas rangés dans une arborescence lisible : ils sont gérés par l'API de fichiers de
Moodle. Chaque fichier est une ligne de la table \texttt{mdl\_files} (préfixe \texttt{mdl\_} par
défaut, voir \texttt{\$CFG->prefix} dans \texttt{config.php}) :

\begin{center}
\small
\begin{tabularx}{\linewidth}{@{}l X@{}}
\toprule
\textbf{Champ} & \textbf{Valeur} \\
\midrule
\texttt{component} & \texttt{mod\_these} \\
\texttt{filearea} & zone (section~\ref{sec:zones}) ; \texttt{sealed} pour les documents scellés \\
\texttt{itemid} & identifiant du dossier, \texttt{mdl\_these\_dossier.id} \\
\texttt{contextid} & contexte de l'activité : ligne de \texttt{mdl\_context} de niveau 70 dont
\texttt{instanceid} est l'identifiant du module de cours \\
\texttt{filepath}, \texttt{filename} & \texttt{/} puis le nom du fichier (\texttt{subject.pdf}\ldots) \\
\texttt{contenthash} & empreinte \textbf{SHA-1} du contenu, qui donne l'emplacement physique \\
\texttt{filesize}, \texttt{timecreated}, \texttt{timemodified} & taille en octets et dates (horodatage Unix) \\
\bottomrule
\end{tabularx}
\end{center}

Le contenu se trouve dans
\texttt{<moodledata>/filedir/\textit{ab}/\textit{cd}/\textit{abcd\ldots}}, où \textit{ab} et
\textit{cd} sont les deux premiers, puis les deux caractères suivants du \texttt{contenthash}, et
\textit{abcd\ldots} le \texttt{contenthash} complet. Deux fichiers identiques partagent le même contenu
physique. Les lignes dont le nom est \og .\fg{} désignent des répertoires et s'ignorent.

% ===========================================================================
\section{Fichiers à verser}\label{sec:zones}

Seules les zones dont le contenu ne change plus sont à verser :

\begin{center}
\small
\begin{tabularx}{\linewidth}{@{}l X X@{}}
\toprule
\textbf{Zone} & \textbf{Contenu} & \textbf{Nom des fichiers} \\
\midrule
\texttt{sealed} & Documents entièrement signés, figés & \texttt{<document>.pdf} (tableau
section~\ref{sec:scellement}) \\
\texttt{pv} & Procès-verbal signé par le jury, scanné & nom d'origine \\
\texttt{pvadjourned} & Procès-verbaux d'ajournement archivés & \texttt{ajournement-AAAA-MM-JJ-<nom d'origine>} \\
\texttt{subjectarchive} & Dépôt de sujet signé, archivé lors d'un changement de sujet &
\texttt{sujet-AAAA-MM-JJ-<n>.pdf} \\
\texttt{signedscan} & Documents signés hors module, scannés &
\texttt{AAAA-MM-JJ-<document>-<signataires>-<nom d'origine>} \\
\texttt{corrections} & Manuscrit corrigé (après corrections demandées par le jury) & nom d'origine \\
\texttt{manuscript} & Manuscrit déposé pour la soutenance, à verser s'il n'y a pas de manuscrit corrigé & nom d'origine \\
\bottomrule
\end{tabularx}
\end{center}

Les autres zones ne sont pas à verser : pièces de travail (\texttt{subjectdocs},
\texttt{defensedocs}, \texttt{directorreport}), documents régénérés à chaque fois
(\texttt{generated}) et images des signatures (\texttt{signature}), déjà incluses dans les PDF scellés.

% ===========================================================================
\section{Métadonnées}\label{sec:metadonnees}

\paragraph{Le dossier} (\texttt{mdl\_these\_dossier}, clé \texttt{id} = \texttt{itemid} du fichier) :

\begin{center}
\small
\begin{tabularx}{\linewidth}{@{}l X@{}}
\toprule
\textbf{Champ} & \textbf{Contenu} \\
\midrule
\texttt{devenumber} & numéro d'enregistrement DEVE (\texttt{2026-042}), attribué à la validation du
dossier de soutenance ; vide avant \\
\texttt{userid} & étudiant (\texttt{mdl\_user} : \texttt{username}, \texttt{lastname},
\texttt{firstname}, \texttt{idnumber}) \\
\texttt{title}, \texttt{titleen} & titre de la thèse, en français et en anglais \\
\texttt{directorid} & directeur de thèse (\texttt{mdl\_user}) \\
\texttt{defensedate} & date de soutenance (horodatage Unix, 0 si non fixée) \\
\texttt{defenseresult} & 0 : pas encore de décision, 1 : admis, 2 : ajourné \\
\texttt{mention} & \texttt{honorable}, \texttt{treshonorable}, \texttt{felicitations} (vide sinon) \\
\texttt{confidential} & 1 : thèse confidentielle (titre et contenu à ne pas divulguer) \\
\texttt{bordereau} & JSON du bordereau ABES ; clés \texttt{embargo} et \texttt{embargountil} pour
une diffusion différée \\
\texttt{status} & étape du circuit : 100 clos, $-1$ abandonné (voir le guide) \\
\bottomrule
\end{tabularx}
\end{center}

\paragraph{L'historique} (\texttt{mdl\_these\_log}) : une ligne par événement, avec
\texttt{dossierid}, \texttt{logaction}, \texttt{remark} et \texttt{timecreated}. Les actions utiles à
l'archivage sont \texttt{documentsealed} (empreinte du scellé), \texttt{signed} (chaque signature) et
\texttt{closed} (clôture du dossier).

% ===========================================================================
\section{Deux façons d'accéder aux fichiers}\label{sec:acces}

\paragraph{Par la base et \texttt{moodledata} (recommandé).} Le connecteur lit
\texttt{mdl\_files} et les tables du module avec un compte de base de données \textbf{en lecture seule},
et lit les fichiers dans \texttt{moodledata/filedir}. C'est la méthode du script fourni
(section~\ref{sec:script}). Exemple de compte PostgreSQL :

\begin{Verbatim}
CREATE ROLE ged LOGIN PASSWORD '...';
GRANT CONNECT ON DATABASE moodle TO ged;
GRANT USAGE ON SCHEMA public TO ged;
GRANT SELECT ON mdl_files, mdl_context, mdl_course_modules, mdl_modules,
                mdl_these_dossier, mdl_these_log, mdl_user TO ged;
\end{Verbatim}

Le compte système qui exécute le connecteur n'a besoin que du droit de lecture sur
\texttt{moodledata/filedir}.

\paragraph{Par le site.} Un compte de service ayant les droits de la DEVE ou de la bibliothèque
télécharge un document de l'établissement par
\nolinkurl{https://<site>/mod/these/document.php?id=<module>&dossierid=<dossier>&type=<document>},
\texttt{<document>} étant la clé du document (\texttt{subject}, \texttt{agreement}\ldots). Si le
document est scellé, c'est le PDF scellé qui est renvoyé ; sinon, le document dans son état de signature
du moment. Les autres zones se téléchargent par
\nolinkurl{https://<site>/pluginfile.php/<contextid>/mod_these/<zone>/<dossier>/<fichier>} (la zone
\texttt{sealed} ne passe pas par \texttt{pluginfile.php}). Cette voie convient à un versement ponctuel ;
pour un versement régulier, préférer la première.

% ===========================================================================
\section{Requêtes}\label{sec:requetes}

\paragraph{Documents scellés, avec leur dossier.} Requête valable pour PostgreSQL et MariaDB/MySQL
(adapter le préfixe \texttt{mdl\_}) :

\begin{Verbatim}
SELECT d.id AS dossierid, d.devenumber, u.lastname, u.firstname, d.title,
       f.filename, f.contenthash, f.filesize, f.timemodified
  FROM mdl_files f
  JOIN mdl_context c ON c.id = f.contextid AND c.contextlevel = 70
  JOIN mdl_course_modules cm ON cm.id = c.instanceid
  JOIN mdl_modules m ON m.id = cm.module AND m.name = 'these'
  JOIN mdl_these_dossier d ON d.id = f.itemid AND d.theseid = cm.instance
  JOIN mdl_user u ON u.id = d.userid
 WHERE f.component = 'mod_these' AND f.filearea = 'sealed'
   AND f.filename <> '.'
 ORDER BY d.id, f.filename;
\end{Verbatim}

\paragraph{Toutes les zones à verser.} Remplacer la condition sur \texttt{filearea} par
\texttt{f.filearea IN ('sealed', 'pv', 'pvadjourned', 'subjectarchive', 'signedscan',
'corrections', 'manuscript')}. Pour un versement incrémental, ajouter
\texttt{AND f.timemodified >= <horodatage du dernier passage>}.

\paragraph{Empreintes SHA-256 des scellés.}

\begin{Verbatim}
SELECT dossierid, remark, timecreated
  FROM mdl_these_log
 WHERE logaction = 'documentsealed'
 ORDER BY dossierid, timecreated;
\end{Verbatim}

La remarque a la forme \og Dépôt de sujet de thèse — SHA-256 f63cad99c472\ldots\fg{} : l'empreinte se
retrouve avec l'expression régulière \verb|SHA-256 ([0-9a-f]{64})|. Un dossier peut avoir plusieurs
lignes pour un même document : une par scellé successif.

% ===========================================================================
\section{Contrôles d'intégrité}\label{sec:integrite}

\begin{enumerate}
  \item \textbf{Fichier physique.} L'empreinte SHA-1 du fichier lu dans \texttt{filedir} doit être égale
        à son \texttt{contenthash}. Sinon, le fichier est absent ou altéré : ne pas le verser, alerter.
  \item \textbf{Document scellé.} L'empreinte SHA-256 du PDF doit figurer parmi celles inscrites à
        l'historique du dossier (\texttt{documentsealed}). C'est la preuve que le fichier versé est
        celui que le module a scellé.
  \item \textbf{Identification d'une version.} Une version se repère par \emph{dossier + document +
        empreinte SHA-256}, jamais par le nom seul (un scellé refait garde le même nom).
\end{enumerate}

Conserver l'empreinte SHA-256 dans les métadonnées de la GED permet de prouver plus tard que le
document n'a pas changé depuis son versement.

% ===========================================================================
\section{Script d'extraction}\label{sec:script}

Le script \texttt{docs/ged/extraction\_ged.py} (Python~3, bibliothèque standard seulement, client
\texttt{psql}) applique les règles ci-dessus pour une base \textbf{PostgreSQL}. Pour MariaDB/MySQL,
seule la fonction \texttt{psql()} est à adapter (client \texttt{mysql}).

\begin{Verbatim}
PGHOST=... PGUSER=ged PGPASSWORD=... PGDATABASE=moodle \
python3 extraction_ged.py --moodledata /var/moodledata --dest /srv/export-ged \
        [--prefix mdl_] [--since 2026-10-01]
\end{Verbatim}

\begin{itemize}
  \item Il ouvre la base en \textbf{lecture seule} (\texttt{default\_transaction\_read\_only}) et ne
        fait que lire \texttt{moodledata}.
  \item Il copie chaque fichier à verser dans
        \texttt{<dest>/<numéro DEVE>/<zone>/<fichier>} (\texttt{dossier-<id>} sans numéro DEVE). Le
        manuscrit initial n'est copié que s'il n'y a pas de manuscrit corrigé.
  \item Il vérifie la SHA-1 de chaque fichier et la SHA-256 de chaque scellé (colonne
        \texttt{scelle\_verifie} : \texttt{oui} ou \texttt{NON}).
  \item Il peut être relancé : un fichier déjà exporté avec la même empreinte est ignoré ; une nouvelle
        version est exportée à côté, sous \texttt{<nom>.<12 premiers caractères de la SHA-256>.pdf},
        jamais par-dessus.
  \item Il complète \texttt{<dest>/manifeste.csv} : dossier, numéro DEVE, étudiant, titre, date de
        soutenance, résultat, mention, zone, fichier, chemin d'export, taille, SHA-256, vérification du
        scellé, date du fichier et date d'export. La GED s'appuie sur ce manifeste pour l'indexation.
  \item Il se termine avec le code 1 si un fichier est absent ou altéré (message sur la sortie
        d'erreur), 0 sinon.
\end{itemize}

\paragraph{Essai.} Sur un Moodle de test : premier passage, un dépôt de sujet scellé exporté et vérifié ;
deuxième passage, rien de nouveau ; scellé levé puis refait, la nouvelle version est exportée à côté de
la première, les deux vérifiées.

% ===========================================================================
\section{Règles pour le connecteur}\label{sec:regles}

\begin{itemize}
  \item \textbf{Lecture seule}, sans exception : ne jamais écrire dans \texttt{filedir} ni dans les
        tables de Moodle, ne jamais supprimer un fichier.
  \item \textbf{Ne jamais supprimer côté GED} une version déjà versée lorsqu'un scellé est refait :
        la nouvelle version s'ajoute.
  \item \textbf{Thèses confidentielles} (\texttt{confidential} = 1) : appliquer dans la GED une
        restriction d'accès équivalente ; le titre ne doit pas être visible de tous.
  \item \textbf{Embargo} (clé \texttt{embargo} du bordereau) : ne pas diffuser le manuscrit avant la
        date \texttt{embargountil}.
  \item \textbf{Données personnelles} : les documents contiennent nom, date et lieu de naissance,
        adresse de l'étudiant et signatures ; les verser sous le régime de conservation prévu pour les
        dossiers d'étudiants.
  \item \textbf{Moment du versement} : de préférence à la clôture du dossier (\texttt{status} = 100),
        quand tous les documents sont scellés ; un versement incrémental (\texttt{--since}) reste sans
        risque grâce à la gestion des versions.
\end{itemize}

% ===========================================================================
\appendix
\section{Script \texttt{extraction\_ged.py}}

\VerbatimInput[fontsize=\scriptsize,numbers=left,numbersep=4pt]{extraction_ged.py}

\end{document}
